\documentclass[10pt]{article}

\usepackage[margin=0.76in]{geometry}
\usepackage{amsmath,amssymb,mathtools}
\usepackage[T1]{fontenc}
\usepackage[utf8]{inputenc}
\usepackage{microtype}
\usepackage{xcolor}
\usepackage{hyperref}
\usepackage{listings}

\hypersetup{
  colorlinks=true,
  linkcolor=blue!45!black,
  urlcolor=blue!45!black,
  citecolor=blue!45!black
}

\lstdefinestyle{pyamplicol}{
  basicstyle=\ttfamily\footnotesize,
  keywordstyle=\bfseries\color{blue!45!black},
  commentstyle=\itshape\color{black!55},
  stringstyle=\color{green!35!black},
  columns=fullflexible,
  keepspaces=true,
  showstringspaces=false,
  frame=single,
  framerule=0.25pt,
  xleftmargin=1.0em,
  xrightmargin=1.0em,
  aboveskip=0.55em,
  belowskip=0.55em
}
\lstset{style=pyamplicol}

\newcommand{\AC}{\textsc{AmpliCol}}
\newcommand{\PAC}{\textsc{pyAmpliCol}}
\newcommand{\J}{\mathcal{J}}
\newcommand{\V}{\mathcal{V}}
\newcommand{\M}{\mathcal{M}}
\newcommand{\heplib}{\texttt{hep\_lib}}

\title{\vspace{-2.0em}Current Recursion and Tensor-Network Execution in \PAC}
\author{}
\date{\vspace{-3.0em}}

\begin{document}
\maketitle
\vspace{-1.2em}

\paragraph{Purpose.}
This note summarizes the matrix-element generation strategy used by \AC\ and
the corresponding \PAC\ implementation.  The common idea is to avoid explicit
Feynman-diagram enumeration.  Instead, matrix elements are assembled from a
directed acyclic graph (DAG) of off-shell currents.  \AC\ emits Fortran code
that numerically sweeps this current table.  \PAC\ builds the same logical DAG
in Python, lowers bounded pieces to tensor-network expressions handled by
Spenso, and compiles the resulting parametric scalar/vector expressions with
Symbolica evaluators.

\section{The AmpliCol Current Recursion}

For a fixed leading-color ordering, a current is labelled by the external
subset it contains, the particle type carried by the off-shell line, and any
spin/chirality information.  Schematically, if $I$ is a set or ordered interval
of external colored legs, a current is computed as
\begin{equation}
  \J^\alpha_{t,\chi}(I)
  =
  P^\alpha{}_\beta(t, p_I)
  \sum_{\substack{I=A\cup B\\A\cap B=\emptyset}}
  \V^\beta{}_{\gamma\delta}(t_A,t_B\to t)\,
  \J^\gamma_{t_A,\chi_A}(A)\,
  \J^\delta_{t_B,\chi_B}(B).
  \label{eq:current-recursion}
\end{equation}
Here $t$ is the current type (gluon, quark, auxiliary tensor, etc.),
$\chi$ is the Weyl chirality where relevant, $p_I=\sum_{i\in I}p_i$, $P$ is
the propagator, and $\V$ is the appropriate Feynman-rule tensor.  The
four-gluon vertex is treated through the auxiliary antisymmetric tensor route,
so it is still represented by cubic current combinations.

\AC's optimized mode does not loop over full helicity configurations as
independent computations.  It weaves helicity information into the current
table itself.  Each external helicity state is a source current with a unique
bit in an \texttt{ext\_cur}-style bitset.  An internal current is identified by
\begin{equation}
  K(\J)=\bigl(t,\chi,\mathrm{bin}(I),\mathrm{ext\_cur}\bigr),
  \label{eq:current-key}
\end{equation}
where $\mathrm{bin}(I)$ records the external subset and $\mathrm{ext\_cur}$ is
the union of the source-current bits feeding the current.  When a new vertex
contribution produces the same key as an existing current, \AC\ appends that
vertex to the existing current rather than creating a duplicate node.  After
all currents are constructed, \texttt{curr2amp} records the two endpoint
currents whose contraction gives each retained helicity amplitude.

The generated Fortran library has the same layered structure:
\begin{lstlisting}[language=Fortran]
subroutine evaluate_amp(p, amps)
  complex(kind=8) :: val_c(1:6, n_cur)
  complex(kind=8) :: int_c(1:6, n_vert)

  call compute_external_currents(pp, val_c)
  do size = 2, n_external_minus_one
    call compute_vertices_size(pp, val_c, int_c)
    call compute_currents_size(pp, val_c, int_c)
  end do
  call compute_amps(amps, val_c)
end subroutine
\end{lstlisting}
The arrays \texttt{val\_c} and \texttt{int\_c} are the numerical current and
interaction tables.  This is the performance-critical idea that \PAC\ keeps:
one topological sweep produces all retained amplitudes, and shared parent
currents are evaluated once.

\section{The pyAmpliCol DAG}

\PAC\ first builds an explicit Python representation of the same current table.
For the leading-color milestone $q\bar q\to Z+n g$, the current identity is
\begin{equation}
  K_{\PAC}=\bigl(\mathrm{pdg},\mathrm{labels},\chi,\mathrm{source\_bits}\bigr).
\end{equation}
The \texttt{source\_bits} field is the direct analogue of \AC's
\texttt{ext\_cur}.  Source currents are created for the allowed quark,
antiquark, gluon, and massive-vector helicity states.  Internal currents are
then generated by increasing current size, with all interactions that feed the
same key attached to one node.

\begin{lstlisting}[language=Python]
for external helicity source s:
    add_current(key=(pdg_s, labels_s, chirality_s, bit_s), is_source=True)

for current_size in increasing_sizes:
    for compatible split A | B:
        result_key = combine_key(key[A], key[B], vertex_kind)
        result = add_or_reuse_current(result_key)
        add_interaction(left=A, right=B, result=result, vertex=vertex_kind)

for retained helicity ancestry:
    add_amplitude(left=current_with_Z_and_gluons, right=anti_current)

prune_dead_currents_after_helicity_filter()
\end{lstlisting}

At runtime, \PAC\ mirrors the Fortran table layout with contiguous arrays:
\[
  C_{s,i,c}\equiv\texttt{current\_rows}[s,i,c],
  \qquad
  T_{s,v,c}\equiv\texttt{interaction\_rows}[s,v,c],
\]
where $s$ is the phase-space sample in the current batch, $i$ is a current id,
$v$ is an interaction id, and $c$ is a component index.  External wavefunctions
fill source-current rows once per sample.  Bounded evaluator blocks then fill
interactions and derived currents stage by stage:
\begin{lstlisting}[language=Python]
fill_source_currents(current_rows, momenta, helicity_sources)

for stage in current_size_layers:
    state = ParamBuilder.pack(current_rows, interaction_rows, momenta)
    interaction_rows[:, stage.vertices] = vertex_eval[stage](state)
    current_rows[:, stage.currents]     = current_eval[stage](state)

amps = amplitude_eval(ParamBuilder.pack(current_rows, momenta))
me2  = normalization * sum(weight[h] * abs(amps[h])**2 for h in retained_h)
\end{lstlisting}
This is why the evaluator-only benchmark is meaningful: the expensive
recursion is inside compiled evaluator calls, while Python only performs source
wavefunction construction, parameter packing, and table transfers.

\section{Spenso Tensor Networks and \texttt{hep\_lib}}

Spenso is used to express Lorentz, spinor, auxiliary-tensor, and eventually
color contractions in a representation-aware way.  A tensor-network expression
contains named tensors with typed slots:
\[
  \V_{\mu\nu}^{\phantom{\mu\nu}\rho}(p_A,p_B)\,
  A^\mu\,B^\nu,
  \qquad
  \text{or in code: }
  \texttt{V(m("a"),m("b"),m("out")) * A(m("a")) * B(m("b"))}.
\]
The expression is not expanded.  It is a compact product of tensor factors with
shared representation labels.  The library \heplib\ is the dictionary that
tells Spenso what each tensor name means: dense model tensors, metric tensors,
momentum-dependent propagators, and parametric rank-one tensors representing
runtime inputs.

\begin{lstlisting}[language=Python]
library = model.build_tensor_library()      # returns TensorLibrary.hep_lib_atom()
builder = ParamBuilder()
mink = Representation.mink(4)

builder.register_rank1_tensor(
    library,
    tensor_name="current_A",
    representation=mink,
    head=("stage", "A"),
    length=4,
    role="block_current",
)
builder.register_rank1_tensor(
    library,
    tensor_name="p_A",
    representation=mink,
    head=("stage", "pA"),
    length=4,
    role="block_momentum",
    real_valued=True,
)

raw = V3g(mink("a"), mink("b"), mink("out"), p_A="p_A", p_B="p_B") \
    * TensorName("current_A")(mink("a")).to_expression() \
    * TensorName("current_B")(mink("b")).to_expression()

network = TensorNetwork(raw, library)
network.execute(library=library)
out_tensor = network.result_tensor(library)
evaluator = out_tensor.evaluator(params=builder.parameter_symbols())
\end{lstlisting}

The important point is that \heplib\ stays fixed during the execution.  New
graph currents are not registered as permanent physics tensors.  Instead, the
current values are runtime parameters for a bounded block.  Spenso executes the
factorized tensor network, contracts the representation indices, and returns a
parametric output tensor.  Symbolica then compiles that parametric tensor to an
evaluator.  For high multiplicity, \PAC\ interleaves this process: after a
current or current-size layer is contracted, its numerical output is written
back into the current table and becomes an input parameter for later blocks.
This avoids one giant scalar expression and keeps the DAG close to \AC's
execution model.

If explicit color tensors are introduced in future full-color modes, color
simplification belongs before Spenso execution:
\[
  \texttt{raw} \leftarrow \texttt{idenso.simplify\_color(raw)} ,
  \qquad
  \text{then build } \texttt{TensorNetwork(raw, hep\_lib)}.
\]
For the current leading-color milestone, color is already reduced to scalar
factors associated with the ordered amplitude.

\section{Example: \texorpdfstring{$d\bar d\to Zgg$}{d dbar -> Z g g}}

The user-level process is written as
\[
  d(p_1)\,\bar d(p_2)\to Z(p_Z)\,g(p_3)\,g(p_4).
\]
For the leading-color current recursion, the colored chain is treated in the
canonical order
\[
  \bar d \;-\; d \;-\; g_3 \;-\; g_4,
\]
with the color-singlet $Z$ inserted on the quark line.  The relevant source
currents are
\[
  \J_{\bar q}(1,h_1),\quad
  \J_q(2,h_2),\quad
  \epsilon^\mu(3,h_3),\quad
  \epsilon^\nu(4,h_4),\quad
  \epsilon_Z^\rho(Z,h_Z).
\]
The gluon subcurrents include
\begin{align}
  \J_g^\alpha(3,4;h_3,h_4)
   &= P_g{}^\alpha{}_\beta(p_3+p_4)
      \V_{3g}^{\beta\mu\nu}(p_3,p_4)
      \epsilon_\mu(3,h_3)\epsilon_\nu(4,h_4),\\
  \J_T^A(3,4;h_3,h_4)
   &= \V_{ggT}^{A\mu\nu}\epsilon_\mu(3,h_3)\epsilon_\nu(4,h_4),
\end{align}
where $\J_T$ is the auxiliary tensor current used to reproduce the four-gluon
structure in larger examples.  The quark line then forms
\begin{align}
  \J_q(2,3) &= P_q(p_2+p_3)\,\V_{qgq}\bigl(\J_q(2),\epsilon(3)\bigr),\\
  \J_q(2,3,4) &=
    P_q(p_2+p_3+p_4)\left[
      \V_{qgq}\bigl(\J_q(2),\J_g(3,4)\bigr)
      + \V_{qgq}\bigl(\J_q(2,3),\epsilon(4)\bigr)
    \right],
\end{align}
and the $Z$ can be inserted before, between, or after the ordered gluons:
\begin{align}
  \J_q(2,Z) &= P_q(p_2+p_Z)\,\V_{qZq}\bigl(\J_q(2),\epsilon_Z\bigr),\\
  \J_q(2,3,Z) &=
    P_q(p_2+p_3+p_Z)\left[
      \V_{qZq}\bigl(\J_q(2,3),\epsilon_Z\bigr)
      + \V_{qgq}\bigl(\J_q(2,Z),\epsilon(3)\bigr)
    \right],
\end{align}
with analogous terms leading to $\J_q(2,3,4,Z)$.  The helicity amplitude is
finally an endpoint contraction,
\begin{equation}
  \M(h_1,h_2,h_3,h_4,h_Z)
  =
  \J_q(2,3,4,Z;h_2,h_3,h_4,h_Z)
  \cdot
  \J_{\bar q}(1,h_1).
\end{equation}
Only retained helicity ancestries survive the filter.  The reported squared
matrix element has the same structure as in \AC:
\begin{equation}
  |\M|^2_{\mathrm{LC}}
  =
  \mathcal{N}_{\mathrm{AC}}
  \sum_{h\in H_{\mathrm{kept}}}
  w_h\,
  \left|\M_h\right|^2,
\end{equation}
where $\mathcal{N}_{\mathrm{AC}}$ contains the leading-color normalization,
couplings, symmetry factors, and initial-state averaging conventions matched
to the Fortran implementation.

\section{Toward a Single Aliased Multi-Output Evaluator}

The current milestone deliberately stays close to the generated Fortran
execution model.  Python owns the current table and performs the topological
sweep at runtime.  Each step of that sweep calls bounded Spenso/Symbolica
evaluator blocks that fill interaction rows, derived-current rows, and finally
the retained amplitude vector.  This is robust and memory-safe: no full
matrix-element expression is ever expanded, and each block has controlled
rank, parameter count, and output count.  The price is that every phase-space
batch crosses the Python boundary several times:
\[
  C_{\mathrm{old}}
  =
  C_{\mathrm{source}}
  \xrightarrow{\text{Python loop}}
  E_1(C,p)
  \xrightarrow{\text{Python table update}}
  E_2(C,p)
  \xrightarrow{\cdots}
  E_{\mathrm{amp}}(C,p).
\]
Here $C$ denotes the current table and each $E_k$ is a compiled evaluator for
one layer or bounded block.  The next natural step is to keep the same
\AC-style DAG construction but move this orchestration from runtime to
generation time.

In that alternative design, the Python graph builder would still enumerate the
same source currents, internal currents, interactions, and endpoint amplitude
records.  However, instead of executing a current block numerically and writing
the result back to \texttt{current\_rows}, every component of every generated
current would receive an opaque symbolic alias:
\[
  J(i,c)
  \quad\longleftrightarrow\quad
  \text{component } c \text{ of current } i.
\]
The right-hand side of this alias is the Spenso-executed parametric tensor
component for that current.  It may contain external wavefunction parameters,
momentum parameters, model constants, and aliases of parent current
components.  The final roots passed to Symbolica are the retained amplitude
components, expressed only in terms of endpoint current aliases and external
parameters:
\[
  \bigl(\M_{h_1},\ldots,\M_{h_N}\bigr)
  =
  R\bigl(J(i_1,c_1),J(i_2,c_2),\ldots;p,\epsilon,u,\bar u\bigr).
\]
The evaluator is then a single multi-output object for the full amplitude
vector.  Runtime evaluation becomes one parameter-pack operation and one
compiled evaluator call, followed by the same helicity/color weighted
$|\M_h|^2$ reduction.

\begin{lstlisting}[language=Python]
aliases = {}
alias_tensors = {}

for current in topological_currents:
    tensor_expr = build_spenso_current_tensor(
        current,
        parent_alias_tensors=alias_tensors,
        hep_lib=hep_lib,
    )
    executed = TensorNetwork(tensor_expr, hep_lib).execute()

    components = {}
    for component, atom in executed.components():
        handle = J(current.id, component)
        aliases[handle] = atom
        components[component] = handle
    alias_tensors[current.id] = tensor_from_alias_components(components)

roots = [
    build_amplitude_component(amplitude, alias_tensors)
    for amplitude in retained_amplitudes
]

evaluator = AliasedExpression(roots, aliases).evaluator(params)
\end{lstlisting}

The important distinction from registering new tensors in \heplib\ is that
\heplib\ remains the fixed dictionary of physics and representation objects:
metrics, propagators, vertices, and parametric external tensors.  Current
aliases are not new model tensors.  They are scalar Symbolica handles for
already-executed tensor-network components.  Spenso is still responsible for
contracting Lorentz, spinor, and auxiliary-tensor indices in each current
definition.  Symbolica is then responsible for compiling the alias graph into
machine code.

The alias mechanism matters.  A naive Symbolica \texttt{functions} map is not
the right representation for the full current pool, because such function
definitions are inlined during evaluator construction and only later rediscovered
by common-subexpression elimination.  That may recreate the large monolithic
expressions that the staged design avoids.  The intended primitive is
\texttt{AliasedAtom}, or a Python binding exposing the same alias-aware
evaluator construction.  Aliases are designed to be opaque handles with a
separate map from handle to body.  This should let the evaluator builder see
the DAG as a pool of reusable current components rather than as a fully
expanded tree.

The generation pass for this backend would therefore have two invariants.
First, the current graph must remain topologically ordered, so every alias body
references only external parameters and aliases that have already been
registered.  Second, each alias handle must be an indeterminate accepted by the
Symbolica evaluator builder, for example a tagged function atom
\texttt{J(i,c)} or a generated symbol \texttt{J\_i\_c}.  The first form is more
diagnostic because it keeps the current id and component visible in printed
debug output.  The second form may be useful if the Python binding exposes only
symbol aliases initially.  In either case, the alias map is acyclic:
\[
  J(i,c) = F_{i,c}\bigl(p,\epsilon,u,\bar u,
    \{J(j,d)\;|\;j\in\mathrm{parents}(i)\}\bigr),
  \qquad
  j < i .
\]

Helicity sharing would be preserved exactly as in the present DAG evaluator.
The source-current ancestry bits are still part of the current key, so two
currents with the same \AC\ identity share one alias tensor, while currents
with different helicity ancestry remain distinct.  The retained amplitude roots
are not constructed by looping over full helicity configurations independently;
they are read from the same filtered current table.  Thus the single-evaluator
backend is not a return to a diagrammatic or per-helicity construction.  It is
the same shared current table, but serialized into an alias graph rather than
into a sequence of runtime stage calls.

The implementation should be validated with explicit size diagnostics before
large multiplicities are attempted.  For every generated process, the artifact
should record
\[
  n_{\mathrm{alias}},\quad
  n_{\mathrm{root}},\quad
  \sum_i \mathrm{ops}(F_i),\quad
  \max_i \mathrm{ops}(F_i),\quad
  \mathrm{bytes}(\mathrm{root},\mathrm{aliases}),
\]
alongside the existing current, interaction, and amplitude counts.  A small
MRE should compare three objects for low multiplicity: the current staged
compiled evaluator, the fully aliased evaluator, and Fortran \AC.  The printed
input to evaluator construction must still show a factorized root plus aliases;
calling an operation that unfolds aliases should be reserved for debugging
only.  If evaluator construction starts from an unfolded expression, this
backend has failed its main design requirement.

There is also a useful intermediate mode.  Instead of putting every current in
one global alias pool immediately, \PAC\ can build one aliased evaluator per
current-size layer, with aliases shared inside the layer and explicit current
table outputs between layers.  This would test the Python binding and alias
semantics on smaller objects while still reducing the number of runtime calls.
Once that is validated, adjacent layers can be fused until the final artifact
contains the full amplitude vector.  The final target remains one evaluator,
but staged fusion gives a controlled path for diagnosing memory and compile
time failures.

For $d\bar d\to Zgg$, the alias pool would start with source handles such as
\[
  J(q_2,h_2,c),\qquad
  J(\bar q_1,h_1,c),\qquad
  J(g_3,h_3,\mu),\qquad
  J(g_4,h_4,\nu),\qquad
  J(Z,h_Z,\rho).
\]
The two-gluon current becomes, component by component,
\begin{equation}
  J(g_{34},\alpha)
  :=
  P_g{}^\alpha{}_\beta(p_3+p_4)\,
  \V_{3g}^{\beta\mu\nu}(p_3,p_4)\,
  J(g_3,\mu)\,J(g_4,\nu),
\end{equation}
and the quark-line current with a $Z$ insertion can reference both source and
derived aliases:
\begin{equation}
  J(q_{234Z},c)
  :=
  \left[
    P_q\,\V_{qgq}\bigl(J(q_2),J(g_{34})\bigr)
    +
    P_q\,\V_{qgq}\bigl(J(q_{23Z}),J(g_4)\bigr)
    + \cdots
  \right]_c .
\end{equation}
The final retained roots are endpoint contractions such as
\begin{equation}
  R_h =
  \sum_c J(q_{234Z};h,c)\,J(\bar q_1;h,c),
\end{equation}
one root for each retained helicity ancestry.  A single evaluator would return
all $R_h$ values at once.  The summed matrix element remains a lightweight
post-processing operation unless the final weighted norm is also added as an
extra root.

This direction has three expected advantages.  First, runtime transfers should
drop substantially because intermediate current tables no longer cross the
Python boundary.  Second, Symbolica can optimize across the complete retained
amplitude vector rather than only within bounded stage blocks.  Third, the
resulting artifact more closely resembles a generated library: load one
evaluator, pack parameters, evaluate all amplitudes.  The main risks are
generation memory, compile time, and alias semantics.  The implementation must
prove that alias definitions remain preserved through evaluator construction
and are not expanded before optimization.  It will likely require adding a
small Python binding for Symbolica's alias-aware evaluator API before this
backend can be tested fairly.

\section{Why This Matches AmpliCol but Remains Python-Friendly}

The Python layer owns graph construction, process metadata, deterministic
phase-space steering, evaluator caching, and validation against \AC.  It does
not perform the hot algebraic recursion component by component.  The hot path
is a sequence of compiled evaluator calls over batches of phase-space points:
\[
  \text{source wavefunctions}
  \;\longrightarrow\;
  \text{stage evaluators}
  \;\longrightarrow\;
  \text{amplitude/raw-sum evaluator}.
\]
Thus \PAC\ keeps \AC's essential optimization---shared currents across
helicity parents and a topological current-table sweep---while replacing
generated Fortran subroutines by Spenso/ Symbolica evaluator artifacts that can
be saved, loaded, and benchmarked independently.

\end{document}
